\documentclass[10pt,letterpaper]{article}
\usepackage{cornell-notes}

\title{Annex C.07.06: 07-06-technology-viewpoint}
\author{}
\date{}
\setDocTitle{Annex C.07.06: 07-06-technology-viewpoint}
\setDocSubtitle{ISO/IEC/IEEE 42010:2022}
\setDocOwner{ISO 42010 Cornell Notes}
\setDocDate{\today}
\setCornellCollection{ISO/IEC/IEEE 42010:2022 Cornell Notes}
\setCornellUnitType{Annex}
\setCornellUnitNumber{C.07.06}
\setCornellUnitTitle{07-06-technology-viewpoint}
\hypersetup{pdftitle={Annex C.07.06: 07-06-technology-viewpoint},pdfsubject={Cornell notes},pdfauthor={}}

\begin{document}
\maketitle
\begin{CornellOverviewBox}{Overview}
{\Large\bfseries\textcolor{CNavy}{Annex C.7.6: Technology viewpoint}}\\[3pt]
{\small\textbf{Classification:} Informative annex}\\
{\small\textbf{Learning objective:} Use the informative guidance on technology viewpoint to reinforce the normative and conceptual material without treating the guidance itself as mandatory.}
\end{CornellOverviewBox}
\vspace{3pt}

\begin{CornellNotesTable}
\CornellNoteRow{Purpose / key concept}{The technology viewpoint frames these concerns: the selection of implementable standards for the system, their implementation and testing.}
\CornellNoteRow{Active recall}{\begin{itemize}[leftmargin=*,nosep,topsep=1pt]
\item What is the central role of Technology viewpoint?
\item How does Technology viewpoint connect to adjacent architecture-description concepts?
\end{itemize}}
\end{CornellNotesTable}


\begin{CornellSummaryBox}{Summary}
The technology viewpoint frames these concerns: the selection of implementable standards for the system, their implementation and testing. This material is informative guidance and does not itself create a conformance requirement.
\end{CornellSummaryBox}

\end{document}
